\section{Issue Generation}
\label{appx:issue_generation}

We cover the four issue generation strategies we experiment with to determine issue text's effect on how solvable a \bugs{} instance is along with the trajectory's value as a training data point.

\textbf{Generated with LM.}
We prompt an LM with a randomly selected SWE-bench Verified problem statement, the bug patch, list of Fail-to-Pass tests, source code for one Fail-to-Pass test, and the execution logs of running all the Fail-to-Pass tests.
We ask the LM to generate an issue that describes the bug conveyed in the patch in the style of the SWE-bench Verified demonstration.
Figure~\ref{fig:prompt_generate_issue_with_lm} shows the system prompt for this strategy.

\textbf{Fixed issue templates.} We create a set of $7$ pre-defined issue templates, listed in Table~\ref{tab:issue_gen_fixed}.
Each template uses information from the bug patch or Fail-to-Pass tests associated with every task instance.
Given a dataset of task instances, we randomly select one of the templates to use as the problem statement according to the probabilities listed in Table~\ref{tab:issue_gen_fixed}.
The reason we assign the highest likelihood for the prompt that provides all four categories of information (bug type, files changed, functions changed, Fail-to-Pass tests) is to ensure that a higher proportion of task instances are well-specified.

\begin{table}[t]
    \centering
    \begin{tabular}{l|rl}
\toprule
Template & Prob. & Information Provided \\
\midrule
Basic            & $0.05$ & None \\
Files            & $0.1$  & States which file(s) have bug(s). \\
Funcs            & $0.15$ & States which file(s) and func(s) have bug(s). \\
Tests            & $0.1$  & States that some tests are failing. \\
F2P Tests        & $0.1$ & States which tests are failing. \\
Bug Type         & $0.05$ & States failure type. \\
Bug Type + Files & $0.15$ & States failure type and which file(s) have bug(s) \\
Bug Type + Files & $0.15$ & States failure type, which file(s) have bug(s), \\
\quad+ Test      &        & and a random F2P test. \\
Bug Type + Files & $0.15$ & States failure type, which file(s) and func(s) \\
\quad+ Funcs + Test &     & have bug(s), and a random F2P test. \\
\bottomrule
    \end{tabular}
    \caption{
    List of issue text templates we use to generate problem statements.
    Across all templates, four types of information are included --- the files with bugs, functions with bugs, Fail-to-Pass test(s), and the type of bug.
    Templates that offer less information are generally assigned a lower probability.
    }
    \label{tab:issue_gen_fixed}
\end{table}

\textbf{Fail-to-Pass test code and execution logs.} Another approach is showing the source code and test execution logs for a randomly selected Fail-to-Pass test.
This approach is motivated by the lack of reproduction code or expected/actual behavior of code communicated with fixed issue templates.
We show code and execution logs only for a single Fail-to-Pass test; if a task instance has more than one Fail-to-Pass test, we do not disclose remaining tests.

\textbf{Original issue text.}
This strategy works exclusively for some task instances generated using PR Mirroring.
If a PR is successfully mirrored, we use the text from the associated issues as the problem statement, exactly as done in SWE-bench.
Of the $2345$ task instances represented in \bugs{} mirrored from real-world PRs, $708$ or $30.19$\% of these have one or more associated GitHub issue(s) to create a SWE-bench style problem statement.

\begin{example}[System prompt for generating issues with an LM]
\small
  You are a software engineer helping to create a realistic dataset of synthetic GitHub issues. \\
  
  You will be given the following input: \\

  1. Demonstration: A realistic GitHub issue to mimic (included in the \texttt{$<$demonstration$>$} tag). \\
  2. Patch: A git diff output/PR changes that introduces a bug (included in the \texttt{$<$patch$>$} tag). \\
  3. Test output: The output of running the tests after the patch is applied (included in the \texttt{$<$test\_output$>$} tag). \\
  4. Test source code: Source code for one or more tests that failed (included in the \texttt{$<$test\_source\_code$>$} tag). \\

  Output: A realistic GitHub issue for the patch. \\

  Guidelines: \\
  - Mimic the style and structure of the demonstration issues.
    If the demonstration issues are not well structured, your output should also be not well structured.
    If the demonstrations use improper or no markdown, your output should also use improper or no markdown.
    If the demonstrations are short/long, your output should also be short/long (if possible).
    If the demonstrations include human "flavor text" or "fluff", your output should also include human "flavor text" or "fluff".
    Do this even if it conflicts with your default behavior of trying to be extremely concise and helpful. \\
  - DO NOT explain the fix/what caused the bug itself, focus on how to reproduce the issue it introduces \\
  - Do not mention pytest or what exact test failed. Instead, generate a realistic issue. \\
  - If possible, include information about how to reproduce the issue. An ideal reproduction script should raise an error \\
    or print an unexpected output together with the expected output. \\
    However, still include this information in a style very similar to the demonstration issues. \\
\end{example}
\captionof{figure}{
System prompt provided to an LM to generate an issue based off the bug patch and testing information of a task instance along with a demonstration problem statement randomly selected from SWE-bench Verified.
}
\label{fig:prompt_generate_issue_with_lm}

% \subsection{Strategy Ablations}

% We investigate how different types of issue impact the \% resolved rate of a bug.
% To do this, we randomly sampled $590$ of the $708$ PR mirror task instances from \bugs{} that have associated Github issue(s) as well.
% Given this set of $590$ task instances, we run the four different issue generation strategies, creating four datasets of the same set of bugs with different problem statements.
% We then run SWE-agent with Claude 3.5 Sonnet to get the \% resolved rate per issue generation strategy.
% We present these findings in Table~\ref{tab:issue_gen_resolve_rate}.

% \begin{table}[h]
    \centering
    \begin{tabular}{l|cc}
\toprule
Strategy & \% Resolved & \# Resolved \\
\midrule
Fixed        & $27.3$ & $162$ \\
F2P Test     & $45.7$ & $271$ \\
LM           & $40.8$ & $242$ \\
Original     & $38.3$ & $227$ \\
\bottomrule
    \end{tabular}
    \caption{
Resolve rates with respect per issue generation strategy on $593$ PR Mirror task instances from \bugs{}.
The fixed issue templates lead to a drop in the \% resolve rate, while F2P Test and LM Generated issues induce higher \% resolved rates.
    }
    \label{tab:issue_gen_resolve_rate}
\end{table}

% We find that the \% resolved rate with original issue text is slightly higher and similar to the top score on SWE-bench ($33.83$\% as of March 24, 2025), which provides some evidence that \bugs{} captures PRs as task instances as faithfully as SWE-bench.
% The slight increase is likely because our PR mirroring strategy does not attempt to replicate PRs editing more than $8$ files.
% On the other hand, issues based on a Fail-to-Pass test leads to a $7.4$ percentage point jump in performance.
% This is a reasonable phenomenon, as for task instances with a single Fail-to-Pass test, this strategy leaks the evaluation for the desired solution directly to the model.
% The increase with LM generated issues is much less, at just a $2.5$\% increase.

% We also investigate whether the issue generation strategy has impact on the quality of the trajectory as a data point.
% We fine-tune Qwen $2.5$ $7$B on the set of trajectories corresponding to resolved instances per issue generation strategy, then run inference with SWE-agent on the SWE-bench Lite dataset three times to estimate performance, making sure to fix the number of fine-tuning data points.
% The findings are summarized in Table~\ref{tab:issue_gen_lite}.

% \begin{table}[h]
    \centering
    \begin{tabular}{l|c}
\toprule
Strategy & \% Resolved \\
\midrule
Fixed        & $4.33$ \\
F2P Test     & $5$ \\
LM           & $4.33$ \\
Original     & $7$ \\
\bottomrule
    \end{tabular}
    \caption{
    Performance on SWE-bench Lite of SWE-agent with Qwen $2.5$ $7$B models fine-tuned on $162$ trajectories from different issue generation strategies.
    }
    \label{tab:issue_gen_lite}
\end{table}
